% !TeX root = ../main.tex
% 5. fejezet — Felhasználókezelés, jogosultságok és adminisztráció
% A ftsrg/thesis-template-latex sablonba \include-dal illesztendő tartalomfájl.
% Hivatkozott források: Backend/HUNting.API/Controllers/AuthController.cs, UsersController.cs,
%   Services/Auth/JwtTokenService.cs, Configuration/JwtOptions.cs, Services/UserService.cs,
%   Models/Users/ApplicationUser.cs, AccountType.cs, Program.cs, Data/SeedData.cs,
%   Frontend/HUNting/src/context/AuthContext.tsx, lib/authStorage.ts, lib/apiClient.ts,
%   services/authApi.ts, pages/Auth/LoginPage.tsx, RegisterPage.tsx, pages/Profile/ProfilePage.tsx,
%   pages/Admin/AdminDashboardPage.tsx, services/adminApi.ts

\chapter{Felhasználókezelés, jogosultságok és adminisztráció}
\label{chap:felhasznalok}

Ez a fejezet a dolgozat funkcióközpontú gerincének első tagja. A felhasználókezelés a HUNting
minden további funkciójának előfeltétele: a hirdetésfeladás, a társasági tagság, a vadásznapló és
a térképes rétegek egyaránt a bejelentkezett felhasználó személyazonosságára és szerepköreire
épülnek. A fejezet a bevezetőben rögzített egységes sablont követi: először a funkció célját és
felhasználói működését mutatom be, ezt követi az adatmodell, a backend- és a frontend-megvalósítás,
végül a teljes adatfolyam, valamint a biztonsági megfontolások és az ismert korlátok összegzése.

\section{A funkció célja és felhasználói működése}
\label{sec:users-cel}

A rendszer két, jogosultságaiban és felhasználói felületében is eltérő fióktípust különböztet meg,
amelyet a backend az \texttt{AccountType} felsorolt típussal ír le
(\texttt{Models/Users/AccountType.cs}): a \emph{Person} típus az egyéni vadászokat, a
\emph{Society} típus a vadásztársaságokat képviseli. Egyéni fiókkal a felhasználó hirdetést adhat
fel, társasági tagságot igényelhet, vadásznaplót vezethet és személyes térképrétegeket kezelhet.
A társasági fiók ezzel szemben a társaság adminisztrációs felületét éri el: tagfelvételi kérelmeket
bírál el, tagokat és bejegyzéseket kezel -- vadásznaplót azonban nem vezethet, és maga nem
jelentkezhet más társaságba.

A felhasználói folyamat a regisztrációval indul. A regisztrációs űrlapon egy kapcsolóval
választható a fióktípus; társasági fiók esetén az űrlap kiegészül a társaság adataival (név,
kapcsolattartási e-mail és telefon, rövid bemutatkozás). Sikeres regisztráció után a rendszer a
felhasználót azonnal be is lépteti: a válasz már tartalmazza a hozzáférési tokent, így külön
bejelentkezésre nincs szükség. A későbbi munkamenetek a bejelentkezési oldalon indulnak, ahol
e-mail-cím és jelszó megadásával igényelhető új token. A bejelentkezett felhasználó a
profiloldalon tekintheti át fiókadatait és szerepköreit, innen érhetők el a fióktípushoz kötött
felületek (egyéni fióknál a „Társaságaim”, társasági fióknál a „Társaság kezelése” oldal), és itt
található a kijelentkezés is. Védett oldal vendégként történő megnyitásakor az alkalmazás a
bejelentkezési oldalra irányít át.

% TODO: SC-12 képernyőkép (figures/sc-12-login.png) — a LoginPage mojibake-javítása UTÁN készüljön!
\begin{figure}[htbp]
	\centering
	\fbox{\parbox[c][5.5cm][c]{0.85\textwidth}{\centering SC-12: A bejelentkezési oldal képernyőképe (elkészítendő)}}
	\caption{A HUNting bejelentkezési oldala}
	\label{fig:login-screen}
\end{figure}

% TODO: SC-13 képernyőkép (figures/sc-13-register.png) — a fióktípus-váltó és a társasági mezők látszódjanak.
\begin{figure}[htbp]
	\centering
	\fbox{\parbox[c][5.5cm][c]{0.85\textwidth}{\centering SC-13: A regisztrációs oldal társasági mezőkkel (elkészítendő)}}
	\caption{Regisztráció társasági fióktípussal, a kiegészítő társasági adatokkal}
	\label{fig:register-screen}
\end{figure}

\section{Adatmodell}
\label{sec:users-adatmodell}

A felhasználókat az ASP.NET Core Identity keretrendszer kezeli, amelyet a projekt egész számú
elsődleges kulcsokkal használ. Az alkalmazás saját felhasználó-entitása, az
\texttt{ApplicationUser} az \texttt{IdentityUser<int>} osztályból származik, és négy
tartományspecifikus mezővel egészíti ki azt: a kötelező, legfeljebb 100 karakteres
\texttt{FullName} névmezővel, az opcionális \texttt{Phone} telefonszámmal, a létrehozás
időbélyegével (\texttt{CreatedAt}), valamint a fióktípust tároló \texttt{AccountType} mezővel
(\texttt{Models/Users/ApplicationUser.cs}). Az entitás navigációs tulajdonságai a felhasználó
kapcsolatait írják le: az \texttt{OwnedSociety} az általa birtokolt társaságra mutat (egy-egy
kapcsolat), a \texttt{Societies} kollekció a tagságait, a \texttt{MembershipRequests} a benyújtott
tagfelvételi kérelmeit, a \texttt{HuntingLogs} pedig a vadásznapló-bejegyzéseit tartalmazza.

Az Identity alapértelmezett \texttt{AspNet*} előtagú tábláit a \texttt{HuntingContext} saját,
beszédesebb nevekre képezi le (\texttt{Users}, \texttt{Roles}, \texttt{UserRoles},
\texttt{UserClaims}, \texttt{UserLogins}, \texttt{RoleClaims}, \texttt{UserTokens}); az átnevezést
a \texttt{RenameIdentityTables} migráció vezette át az adatbázissémán. Az e-mail-cím egyedisége
két szinten is biztosított: az Identity konfigurációjában a \texttt{RequireUniqueEmail} beállítás,
az adatbázisban pedig egyedi index garantálja. A társaság és tulajdonosa közötti egy-egy
kapcsolatot a \texttt{Society.OwnerUserId} idegen kulcs valósítja meg; a tulajdonos felhasználó
törlésekor a kapcsolat \texttt{SetNull} viselkedéssel szakad el, így a társaság rekordja nem vész
el. Megjegyzendő, hogy a \texttt{Society} entitás saját \texttt{Email} és \texttt{PasswordHash}
mezőkkel is rendelkezik, a bejelentkezés azonban kizárólag a felhasználói fiókon keresztül
történik: társasági regisztrációkor a \texttt{PasswordHash} mezőbe csupán egy véletlen generált
helykitöltő érték kerül.

Az alkalmazásszintű szerepköröket az Identity szereptáblája tárolja. A kezdeti adatfeltöltés
(\texttt{Data/SeedData.cs}) három szerepet hoz létre: az \emph{Administrator} a teljes rendszert
felügyeli, a \emph{SocietyManager} a társasági fiókok alapértelmezett szerepe, a \emph{Hunter}
pedig az egyéni vadászoké. Fontos megkülönböztetni ezektől a társaságon belüli szerepeket
(\texttt{SocietyRole}: Vezeto, Gazda, Vador, Tag), amelyek nem az Identity részei, hanem a tagsági
kapcsolatrekordhoz kötődnek; ezeket a~\ref{chap:tarsasagok}.~fejezet tárgyalja.

\section{Backend-megvalósítás}
\label{sec:users-backend}

A hitelesítési funkciókat az \texttt{AuthController} valósítja meg, amely -- a rétegzett
felépítéstől némileg eltérve -- nem egy dedikált szolgáltatásosztályra, hanem közvetlenül az
Identity \texttt{UserManager} és \texttt{SignInManager} komponenseire, a tokengeneráláshoz pedig
az \texttt{ITokenService} absztrakcióra támaszkodik. A~\ref{tab:auth-endpoints}.~táblázat a
kontroller végpontjait foglalja össze.

\begin{table}[htbp]
	\centering
	\caption{Az \texttt{AuthController} végpontjai}
	\label{tab:auth-endpoints}
	\begin{tabular}{|l|l|p{6.6cm}|}
		\hline
		\textbf{Végpont} & \textbf{Hozzáférés} & \textbf{Feladat} \\
		\hline
		POST \texttt{/api/Auth/register} & anonim & Fiók létrehozása (Person vagy Society), azonnali tokenkibocsátással \\
		\hline
		POST \texttt{/api/Auth/login} & anonim & Hitelesítés, JWT és felhasználói adatok visszaadása \\
		\hline
		POST \texttt{/api/Auth/logout} & bejelentkezett & Kijelentkezés (állapotmentes, 204-es válasz) \\
		\hline
		GET \texttt{/api/Auth/me} & bejelentkezett & Az aktuális felhasználó profilja, szerepei és birtokolt társasága \\
		\hline
	\end{tabular}
\end{table}

\subsection{Regisztráció}

A \texttt{Register} akció a kérés DTO-jából (\texttt{RegisterRequest} rekord) felépíti az
\texttt{ApplicationUser} példányt, majd a \texttt{UserManager.CreateAsync} hívással hozza létre a
fiókot; a jelszó hash-eléséről az Identity gondoskodik. Sikertelen létrehozás esetén a végpont
szándékosan általános hibaüzenettel tér vissza, a konkrét okok felsorolása nélkül. Ezt követően a
fiók a típusának megfelelő alapértelmezett szerepet kapja: társasági fiók esetén
\emph{SocietyManager}, egyébként \emph{Hunter} szerepet.

Társasági regisztrációnál a folyamat két további lépéssel bővül. A kontroller ellenőrzi, hogy a
kérés tartalmazza-e a társaság nevét -- ennek hiányában hibával tér vissza --, majd létrehozza a
\texttt{Society} entitást, tulajdonosként az imént regisztrált felhasználót megjelölve. Ha a
seed-adatok között létezik a „Vezeto” elnevezésű társasági szerep, a rendszer a tulajdonost azonnal
taggá is teszi ezzel a szereppel, így a társasági fiók a regisztráció pillanatától teljes körűen
kezelheti saját társaságát. A válasz mindkét ágon a bejelentkezésével azonos szerkezetű: a
kibocsátott tokent, annak lejáratát és a felhasználó adatait (\texttt{UserDto}) tartalmazza.

\subsection{Bejelentkezés és tokenkibocsátás}

A \texttt{Login} akció először e-mail-cím alapján keresi meg a felhasználót
(\texttt{FindByEmailAsync}), majd a \texttt{SignInManager.CheckPasswordSignInAsync} hívással
ellenőrzi a jelszót. Mindkét hibaág ugyanazt az általános üzenetet adja vissza („Hibás e-mail cím
vagy jelszó.”), így a válaszból nem derül ki, hogy létezik-e fiók az adott címmel -- ez a
felhasználó-felsorolásos (user enumeration) támadások elleni bevett védekezés. Sikeres hitelesítés
után a \texttt{JwtTokenService} állítja elő a hozzáférési tokent
(\ref{lst:jwt-token}.~kódrészlet).

\begin{lstlisting}[language={[Sharp]C}, caption={A JWT előállítása a \texttt{JwtTokenService} osztályban (kivonat)}, label={lst:jwt-token}]
var claims = new List<Claim>
{
    new(JwtRegisteredClaimNames.Sub, user.Id.ToString()),
    new(JwtRegisteredClaimNames.Email, user.Email ?? string.Empty),
    new(JwtRegisteredClaimNames.Jti, Guid.NewGuid().ToString()),
    new(ClaimTypes.NameIdentifier, user.Id.ToString()),
    new(ClaimTypes.Name, user.FullName),
    new(ClaimTypes.Email, user.Email ?? string.Empty)
};

foreach (var role in roles)
{
    claims.Add(new Claim(ClaimTypes.Role, role));
}

claims.Add(new("account_type", user.AccountType.ToString()));

var expires = DateTime.UtcNow.AddMinutes(_options.AccessTokenLifetimeMinutes);

var jwt = new JwtSecurityToken(
    issuer: _options.Issuer,
    audience: _options.Audience,
    claims: claims,
    expires: expires,
    signingCredentials: credentials);
\end{lstlisting}

A token a felhasználó azonosítóját, nevét és e-mail-címét, az összes alkalmazásszintű szerepét,
valamint egy egyedi tokenazonosítót (\texttt{jti}) és a fióktípust leíró saját
\texttt{account\_type} claimet hordozza. Az aláírás szimmetrikus kulcsú HMAC-SHA256 algoritmussal
készül; a kulcsot, a kibocsátót, a célközönséget és az élettartamot (alapértelmezetten 60 perc) a
konfiguráció \texttt{Jwt} szakaszából kötött \texttt{JwtOptions} beállításosztály szolgáltatja. A
bejövő kérések tokenjeit a~\ref{chap:alapok}.~fejezetben bemutatott bearer-middleware validálja
(kibocsátó-, célközönség- és aláírásellenőrzés, egyperces megengedett órakülönbséggel). A
\texttt{Me} végpont a validált tokenből a \texttt{NameIdentifier} claim alapján azonosítja a
felhasználót, és a friss adatbázis-állapotot adja vissza. A \texttt{Logout} végpont állapotmentes:
kizárólag 204-es választ ad, a token érvényességét nem vonja vissza -- ennek következményeit
a~\ref{sec:users-biztonsag}.~alfejezet tárgyalja.

\subsection{A jogosultsági modell}

A hozzáférés-szabályozás két szinten valósul meg. Az első szint deklaratív: a végpontokat az
ASP.NET Core \texttt{[Authorize]} attribútuma védi, amely szerepkör-megkötéssel is ellátható -- a
felhasználó-adminisztrációs kontroller például teljes egészében az
\texttt{[Authorize(Roles = "Administrator")]} attribútummal zárja ki a nem adminisztrátor
hívókat. A második szint imperatív, erőforrás-szintű ellenőrzés: a módosító műveletek a kérdéses
rekord tulajdonosát vetik össze a tokenből kiolvasott felhasználó-azonosítóval, és az
adminisztrátori szerep ez alól kivételt képez. Ezt a mintát szemlélteti
a~\ref{lst:owner-check}.~kódrészlet a hirdetésmódosítás példáján; ugyanez az ellenőrzés ismétlődik
a hirdetéstörlésben, a képkezelő végpontokban, a vadásznaplóban és a társasági bejegyzéseknél is.

\begin{lstlisting}[language={[Sharp]C}, caption={Erőforrás-tulajdonos ellenőrzés az \texttt{AdsController.Update} akcióban (kivonat)}, label={lst:owner-check}]
var existing = await _svc.GetByIdAsync(id);
if (existing == null) return NotFound();

var userId = GetUserId();
if (userId == null) return Unauthorized();

var isAdmin = User.IsInRole("Administrator");
if (!isAdmin && existing.OwnerUserId != userId.Value)
{
    return Forbid();
}

// ... mezok masolasa ...

if (isAdmin)
{
    existing.OwnerUserId = ad.OwnerUserId;
    existing.OwnerSocietyId = ad.OwnerSocietyId;
}
\end{lstlisting}

A jogosultsági modell harmadik eleme a fióktípus szerinti korlátozás: egyes végpontok a token
\texttt{account\_type} claimje, illetve a felhasználó rekordja alapján zárják ki a társasági
fiókokat -- társasági fiók nem hozhat létre vadásznapló-bejegyzést
(\ref{chap:naplo}.~fejezet), és nem nyújthat be tagfelvételi kérelmet
(\ref{chap:tarsasagok}.~fejezet).

\section{Frontend-megvalósítás}
\label{sec:users-frontend}

A kliensoldali hitelesítés központi eleme az \texttt{AuthContext} React-kontextus
(\texttt{src/context/AuthContext.tsx}), amely az egész alkalmazást átfogó szolgáltatóként a
bejelentkezett felhasználót és a hitelesítési műveleteket (\texttt{login}, \texttt{register},
\texttt{logout}, \texttt{refreshProfile}) teszi elérhetővé. A munkamenet két értéken alapul,
amelyeket az \texttt{authStorage} segédmodul a böngésző \texttt{localStorage} tárolójában kezel: a
\texttt{hunting.accessToken} kulcs a JWT-t, a \texttt{hunting.user} kulcs a felhasználó
szerializált profilját tárolja. A kontextus kezdeti állapotát a tárolt profilból állítja helyre,
így az oldal újratöltése nem szakítja meg a munkamenetet; ha tokent talál profil nélkül, a
felcsatolódáskor a \texttt{/auth/me} végpontról frissíti azt, sikertelenség esetén pedig kiüríti a
tárolót.

A bejelentkezési és regisztrációs művelet azonos sémát követ: az \texttt{authApi} szolgáltatásmodul
meghívja a megfelelő végpontot -- a regisztrációnál a kliensoldali \texttt{name} mezőt a backend
által várt \texttt{fullName} mezőre képezve át --, majd sikeres válasz esetén a token és a profil a
tárolóba és a kontextus állapotába kerül. A hibákat a hívó űrlap kapja meg és jeleníti meg. A
kijelentkezés először a backend \texttt{logout} végpontját hívja (ha van tárolt token), majd minden
esetben kiüríti a tárolót és a kontextust.

A kéréseket az \texttt{apiClient} axios-példány két elfogója (interceptor) egészíti ki. A
kérés-elfogó minden hívás \texttt{Authorization} fejlécébe befűzi a tárolt tokent
\texttt{Bearer} sémával. A válasz-elfogó a 401-es státuszú válaszokra globálisan reagál: törli a
hitelesítési kulcsokat és a bejelentkezési oldalra irányítja a böngészőt -- így a lejárt vagy
érvénytelen tokennel dolgozó munkamenet bármely védett hívás hatására automatikusan és egységesen
kijelentkezik.

Az oldalszintű hozzáférés-védelem nem központi útvonalőrrel (route guard), hanem az egyes
oldalkomponensekben valósul meg. A profiloldal vendég látogatóját a \texttt{Navigate} komponens a
bejelentkezési oldalra irányítja; a bejelentkezési és regisztrációs oldal a már hitelesített
felhasználót az útvonalállapotban (\texttt{location.state.from}) átadott eredeti céloldalra --
ennek hiányában a főoldalra -- lépteti tovább; az admin vezérlőpult pedig a szerepkörlista alapján
dönt. A regisztrációs űrlap kliensoldali validációt is végez: a kötelező mezőket és az e-mail
formátumát natív HTML-attribútumok ellenőrzik, a jelszóra hatkarakteres minimum vonatkozik
(összhangban a szerveroldali Identity-szabályzattal), társasági fióktípusnál pedig a beküldés
előfeltétele a társaságnév kitöltése.

% TODO: SC-16 képernyőkép (figures/sc-16-profile.png)
\begin{figure}[htbp]
	\centering
	\fbox{\parbox[c][5.5cm][c]{0.85\textwidth}{\centering SC-16: A profiloldal képernyőképe (elkészítendő)}}
	\caption{A profiloldal a fiókadatokkal, szerepkörökkel és a fióktípusfüggő gyorslinkekkel}
	\label{fig:profile-screen}
\end{figure}

\section{Adminisztrációs felület}
\label{sec:users-admin}

A rendszergazdai funkciók két rétegből állnak. A backendoldali \texttt{UsersController} teljes
felhasználó-adminisztrációt kínál (listázás, lekérdezés, létrehozás, módosítás, törlés), és --
mint fentebb szerepelt -- kizárólag \emph{Administrator} szereppel érhető el. A létrehozó és
módosító végpontok a szerepkiosztást is kezelik: a \texttt{UserService} a még nem létező szerepet
automatikusan megteremti, a módosító akció pedig halmazkülönbség-képzéssel számítja ki az
eltávolítandó és hozzáadandó szerepeket, így a kérésben megadott lista deklaratívan írja le a
kívánt végállapotot.

A frontendoldali admin vezérlőpult (\texttt{AdminDashboardPage}) két nézet között vált:
hirdetések és felhasználók. Mindkét nézet lusta lekérdezést használ -- a React Query
\texttt{enabled} kapcsolójával csak az aktív nézet adatai töltődnek be --, a törlőműveletek pedig
natív megerősítő kérdés (\texttt{window.confirm}) után futnak le, majd a találati lista
gyorsítótárát érvénytelenítve azonnal frissítik a felületet. A felület jelenleg listázásra és
törlésre szorítkozik; a backend létrehozó és módosító végpontjaihoz nem tartozik felületi elem,
ami a~\ref{sec:users-biztonsag}.~alfejezetben felsorolt korlátok egyike. Az adminisztrátor ezen
felül a vadásznaplók teljes állományát is lekérdezheti a \texttt{scope=all} paraméterrel
(\ref{chap:naplo}.~fejezet).

% TODO: SC-14 képernyőkép (figures/sc-14-admin.png) — felhasználók nézet, törlés-megerősítéssel.
\begin{figure}[htbp]
	\centering
	\fbox{\parbox[c][5.5cm][c]{0.85\textwidth}{\centering SC-14: Az admin vezérlőpult képernyőképe (elkészítendő)}}
	\caption{Az adminisztrációs vezérlőpult felhasználókezelő nézete}
	\label{fig:admin-screen}
\end{figure}

\section{Adatfolyam, biztonság és korlátok}
\label{sec:users-biztonsag}

A bejelentkezés teljes adatfolyama a következő lépésekből áll (a folyamatot
a~\ref{fig:auth-sequence}.~ábra szemlélteti):

\begin{enumerate}
	\item A felhasználó kitölti a bejelentkezési űrlapot; a \texttt{LoginPage} az
	      \texttt{AuthContext.login} műveletet hívja.
	\item Az \texttt{authApi} a \texttt{POST /api/auth/login} végpontra küldi az e-mail-címet és a
	      jelszót; a kérés-elfogó ekkor még nem fűz be tokent.
	\item Az \texttt{AuthController} a \texttt{UserManager} segítségével megkeresi a felhasználót,
	      majd a \texttt{SignInManager} ellenőrzi a jelszót a tárolt hash-sel szemben.
	\item Siker esetén a \texttt{JwtTokenService} előállítja az aláírt tokent a szerepkör- és
	      identitás-claimekkel.
	\item A válaszból a kliens a tokent és a profilt a \texttt{localStorage}-ba menti, és frissíti a
	      kontextus állapotát; a felület a bejelentkezett nézetre vált.
	\item A további API-hívások \texttt{Authorization: Bearer} fejléccel utaznak; a middleware
	      kérésenként validálja a tokent, és a claimekből állítja elő a felhasználó identitását.
	\item Érvénytelen vagy lejárt token esetén a backend 401-gyel válaszol, amire a válasz-elfogó a
	      tároló kiürítésével és a bejelentkezési oldalra történő átirányítással reagál.
\end{enumerate}

% TODO: B-01 szekvenciadiagram (figures/b-01-login-sequence.pdf) — LoginPage → AuthContext → axios →
%       AuthController → SignInManager → JwtTokenService → localStorage.
\begin{figure}[htbp]
	\centering
	\fbox{\parbox[c][6.5cm][c]{0.9\textwidth}{\centering B-01: A bejelentkezés és tokenkibocsátás szekvenciadiagramja (elkészítendő)}}
	\caption{A bejelentkezési folyamat szekvenciadiagramja}
	\label{fig:auth-sequence}
\end{figure}

% TODO: B-02 folyamatábra (figures/b-02-request-auth.pdf) — kérések hitelesítése: Bearer-fejléc →
%       validálás → claims → [Authorize]/szerep → erőforrás-tulajdonos ellenőrzés; 401-interceptor ág.
\begin{figure}[htbp]
	\centering
	\fbox{\parbox[c][6.5cm][c]{0.9\textwidth}{\centering B-02: A kérések hitelesítésének folyamatábrája (elkészítendő)}}
	\caption{A védett kérések kiszolgálásának folyamata a tokenvalidálástól az erőforrás-szintű ellenőrzésig}
	\label{fig:auth-flow}
\end{figure}

A megoldás több ponton követi a bevett biztonsági gyakorlatot: a jelszavakat az Identity
alapértelmezett, PBKDF2 alapú hash-elése védi; a bejelentkezési hibaüzenet nem árulja el a fiók
létezését; az e-mail-cím egyedisége adatbázisszinten kikényszerített; a token rövid élettartamú és
aláírása kérésenként ellenőrzött; az erőforrás-szintű tulajdonos-ellenőrzés pedig megakadályozza,
hogy egy hitelesített felhasználó más rekordjait módosítsa.

Ugyanakkor a jelenlegi megvalósítás korlátait is rögzíteni kell. A jelszószabályzat tudatosan
enyhített -- hat karakteres minimum, összetettségi követelmények nélkül --, ami éles üzemben
szigorítandó. A kijelentkezés kizárólag kliensoldali: a kibocsátott token a lejáratáig érvényes
marad, mivel sem visszavonási lista, sem frissítőtoken-mechanizmus nincs. A token
\texttt{localStorage}-ban tárolása egyszerű és újratöltésálló, de XSS-támadás esetén a token
kiolvasható -- ezt a kockázatot a keretrendszer kimeneti kódolása mérsékli, meg nem szünteti. A
társasági regisztráció nem atomi művelet: ha a felhasználó létrehozása után a társaságnév
hiányzik, a kérés hibával tér vissza, miközben a felhasználói fiók már létrejött. Az admin
felhasználó-létrehozó és -módosító végpontok kivételei kezeletlenül 500-as válasszá alakulnak,
felületi funkció pedig nem tartozik hozzájuk. Végül a kezdeti adatfeltöltésben rögzített
adminisztrátori jelszó kizárólag fejlesztési környezetben elfogadható; éles telepítés előtt
cserélendő, ahogyan a konfigurációban tárolt JWT-aláírókulcs is titokkezelőbe való.
